\documentclass[11pt,letterpaper]{article}
\usepackage[utf8]{inputenc}
\usepackage[T1]{fontenc}
\usepackage[spanish,es-nodecimaldot,es-noshorthands,es-tabla]{babel}
\usepackage{lmodern}
\usepackage[margin=2.1cm]{geometry}
\usepackage{amsmath,amssymb}
\usepackage{array,booktabs,multirow}
\usepackage{enumitem}
\usepackage{xcolor}
\usepackage{tikz}
\usepackage[most]{tcolorbox}
\usepackage{fancyhdr}
\usepackage[hidelinks]{hyperref}

\definecolor{azul}{HTML}{1F4E79}
\definecolor{verde}{HTML}{2E6B3E}
\definecolor{rojo}{HTML}{9C2B2B}
\definecolor{naranja}{HTML}{B5651D}
\definecolor{gris}{HTML}{555555}

\newtcolorbox{teoria}[1]{enhanced,breakable,colback=azul!4,colframe=azul,
  title={#1},fonttitle=\bfseries,left=6pt,right=6pt,top=4pt,bottom=4pt}
\newtcolorbox{receta}[1]{enhanced,breakable,colback=verde!5,colframe=verde,
  title={Receta: #1},fonttitle=\bfseries,left=6pt,right=6pt,top=4pt,bottom=4pt}
\newtcolorbox{ejemplo}[1]{enhanced,breakable,colback=white,colframe=gris,
  title={#1},fonttitle=\bfseries,left=6pt,right=6pt,top=4pt,bottom=4pt}
\newtcolorbox{trampa}{enhanced,breakable,colback=rojo!5,colframe=rojo,
  title={Trampa t\'ipica},fonttitle=\bfseries,left=6pt,right=6pt,top=4pt,bottom=4pt}
\newtcolorbox{mecaniza}{enhanced,breakable,colback=naranja!6,colframe=naranja,
  title={Mecaniza (hazlo sin ver y compara)},fonttitle=\bfseries,left=6pt,right=6pt,top=4pt,bottom=4pt}

\newcommand{\E}{\operatorname{E}}
\newcommand{\Var}{\operatorname{Var}}
\newcommand{\Cov}{\operatorname{Cov}}
\newcommand{\fxy}{f_{XY}}
\newcommand{\resp}[1]{\textcolor{verde}{\ensuremath{\boldsymbol{#1}}}}
\newcommand{\origen}[1]{\hfill{\small\normalfont\itshape #1}}

\setlist{itemsep=2pt,topsep=3pt}
\setlength{\parindent}{0pt}
\setlength{\parskip}{5pt}
\pagestyle{fancy}\fancyhf{}
\lhead{\small Probabilidad --- Tema 3}\rhead{\small Variables aleatorias conjuntas}
\cfoot{\thepage}

\begin{document}

\begin{center}
{\LARGE\bfseries Gu\'ia para el examen del Tema 3}\\[4pt]
{\Large Variables aleatorias conjuntas}\\[4pt]
{\small Probabilidad, Facultad de Ingenier\'ia --- basada en tus apuntes, tu formulario y la Serie 3 (2027-1)}
\end{center}

\tableofcontents
\newpage

\section{Radiograf\'ia del examen: qu\'e va a venir}

\subsection{Qu\'e nos dice tu examen del Tema 2}

Tu segundo parcial tuvo \textbf{5 problemas}, cada uno de un subtema distinto, todos de procedimiento
(nada de demostraciones) y con el mismo molde que la serie y los ejemplos de los apuntes:

\begin{center}\small
\begin{tabular}{clp{7.2cm}}
\toprule
\# & Qu\'e pidi\'o & De d\'onde sale el molde\\
\midrule
1 & Dada $F_X$: media, variancia, probabilidad & Serie 2, problema 9 (misma idea, otra funci\'on)\\
2 & Dada $F$ discreta: funci\'on de masa y $P(4<X\le 7)$ & Serie 2, problema 7, \textbf{casi id\'entico}\\
3 & Densidad: media, mediana, moda, sesgo & Serie 2, problemas 13 y 15\\
4 & Hallar $c$, obtener $F_Y$, graficar, usar $F_Y$ & Serie 2, problemas 8 y 13\\
5 & Valor esperado de un costo $120X+10$ & Serie 2, problema 14 / propiedades de $\E$\\
\bottomrule
\end{tabular}
\end{center}

\textbf{Conclusi\'on:} la profesora toma un problema de la serie casi tal cual y arma los dem\'as con el
mismo molde de la serie y de los ejemplos de clase, cambiando la funci\'on o los n\'umeros. Recorre el temario
en orden y pide incisos cortos y encadenados. Si dominas los moldes de la Serie 3 y de los ejemplos 3.3 a 3.13,
no hay sorpresa posible.

\subsection{Pron\'ostico para el examen del Tema 3}

Las probabilidades son \emph{mi estimaci\'on}, hecha con el patr\'on anterior, el peso de cada subtema en tus
apuntes y cu\'antas veces aparece en la serie. No son un dato oficial.

\begin{center}\small
\begin{tabular}{p{6.6cm}cp{5.6cm}}
\toprule
Tipo de ejercicio & Prob. & Molde para practicar\\
\midrule
Tabla conjunta discreta: probabilidades de eventos, marginales, condicional & 95\,\% & Ej.~3.3, Serie 1 y 8\\
Densidad conjunta: hallar $k$ y una probabilidad (regi\'on rectangular o triangular) & 95\,\% & Ej.~3.4, Serie 2, 3 y 7\\
Marginales y condicionales continuas (y probabilidad condicional) & 85\,\% & Ej.~3.6, Serie 3 y 6\\
Covariancia y coeficiente de correlaci\'on & 85\,\% & Serie 9 y 10, Autoexamen 7 y 10\\
?`Son independientes? (como inciso) & 75\,\% & Ej.~3.8, Serie 7e y 8c\\
$\E$ y $\Var$ de una combinaci\'on lineal & 70\,\% & Serie 11, Ej.~3.13, Autoexamen 8 y 9\\
Valor esperado de $g(X,Y)$ con contexto (costo, sobrante) & 65\,\% & Ej.~3.9, Serie 5 y 8d--e\\
Curva de regresi\'on $\E[Y\mid x]$ & 45\,\% & Ej.~3.10, Serie 6c\\
Construir la tabla conjunta desde un experimento & 30\,\% & Ej.~3.1, Serie 10a\\
Funci\'on de distribuci\'on conjunta (tabla acumulada) & 25\,\% & Ej.~3.7\\
Generadora de momentos de una suma & 10\,\% & Ej.~3.11\\
\bottomrule
\end{tabular}
\end{center}

\textbf{Examen m\'as probable (5 problemas):} (1) una tabla discreta con marginales, condicional e independencia;
(2) una densidad con $k$ y una probabilidad sobre un tri\'angulo; (3) marginal y condicional de una densidad;
(4) covariancia y correlaci\'on; (5) combinaci\'on lineal o valor esperado con contexto.

\subsection{C\'omo usar esta gu\'ia}

\begin{enumerate}
\item Lee la secci\'on~\ref{sec:ideas} una vez: son las \'unicas tres ideas del tema.
\item En cada secci\'on, lee la \textbf{receta} (verde), sigue el \textbf{ejemplo resuelto} (gris) con l\'apiz en mano
y despu\'es resuelve el cuadro \textbf{Mecaniza} (naranja) tapando la respuesta.
\item Los cuadros rojos son los errores que m\'as puntos quitan.
\item Al final hay un simulacro de 5 problemas con soluci\'on completa. Hazlo con reloj: 90 minutos.
\end{enumerate}
Todos los resultados de la gu\'ia est\'an verificados y coinciden con las respuestas que trae tu serie.

\section{Las tres ideas del tema (desde cero)}\label{sec:ideas}

En el Tema 2 ten\'ias \emph{una} variable y su funci\'on $f_X(x)$. Ahora el mismo experimento contesta
\emph{dos} preguntas a la vez: $X$ y $Y$. La funci\'on conjunta $\fxy(x,y)$ dice qu\'e tan probable es cada
\emph{pareja} $(x,y)$.

\begin{teoria}{Diccionario discreto $\leftrightarrow$ continuo}
\begin{center}
\begin{tabular}{lcc}
\toprule
 & Discreto (tabla) & Continuo (f\'ormula)\\
\midrule
$\fxy(x,y)$ es & $P(X=x,\,Y=y)$ & altura de una superficie\\
Todo suma & $\sum\sum \fxy = 1$ & $\iint \fxy\,dx\,dy = 1$\\
Probabilidad de un evento & sumar las celdas del evento & integrar sobre la regi\'on\\
\bottomrule
\end{tabular}
\end{center}
Todo lo que se hace con sumas en una tabla se hace con integrales en una densidad. Es la misma receta.
\end{teoria}

\begin{teoria}{Idea 1 --- Marginal: \guillemotleft olvidar\guillemotright{} una variable}
Si s\'olo me interesa $X$, junto todo lo que pasa con $Y$:
\[
f_X(x)=\sum_{y}\fxy(x,y)\qquad\text{o}\qquad f_X(x)=\int_{-\infty}^{\infty}\fxy(x,y)\,dy .
\]
Se suma o se integra \textbf{sobre la variable que quiero desaparecer}. El resultado s\'olo depende de $x$.
En una tabla son los totales de rengl\'on y de columna (por eso se llaman \emph{marginales}: van al margen).
\end{teoria}

\begin{teoria}{Idea 2 --- Condicional: \guillemotleft ya s\'e cu\'anto vali\'o la otra\guillemotright}
\[
f_{X\mid y_0}(x\mid Y=y_0)=\frac{\fxy(x,y_0)}{f_Y(y_0)}
\qquad\qquad
f_{Y\mid x_0}(y\mid X=x_0)=\frac{\fxy(x_0,y)}{f_X(x_0)}
\]
\textbf{Conjunta entre la marginal de lo que ya conozco.} En una tabla: tomo el rengl\'on (o columna) del valor
conocido y lo divido entre su total, para que vuelva a sumar 1.
\end{teoria}

\begin{teoria}{Idea 3 --- Independencia: una variable no informa sobre la otra}
\[
X,\,Y \text{ independientes} \iff \fxy(x,y)=f_X(x)\,f_Y(y)\quad\text{para \textbf{todo} } x,\,y .
\]
Si falla en una sola pareja, son dependientes. Si son independientes entonces
$\E(XY)=\E(X)\E(Y)$, $\Cov(X,Y)=0$ y $\rho=0$.
\end{teoria}

Con estas tres ideas m\'as el valor esperado
$\E[g(X,Y)]=\sum\sum g\,\fxy$ (o $\iint g\,\fxy$) se resuelve todo el examen.

\section{Variables conjuntas discretas (la tabla)}

\begin{receta}{cualquier problema con tabla}
\begin{enumerate}
\item \textbf{Ubica qui\'en es qui\'en:} qu\'e variable va en renglones y cu\'al en columnas. Escr\'ibelo.
\item \textbf{Agrega los m\'argenes:} total de cada rengl\'on y de cada columna. Verifica que todo sume 1.
\item \textbf{Probabilidad de un evento} ($X>Y$, $X+Y=4$, $X=Y$\ldots): marca las celdas que cumplen y s\'umalas.
\item \textbf{Condicional:} toma el rengl\'on o columna del valor dado y divide cada celda entre su total.
\item \textbf{Independencia:} compara cada celda con (total de su rengl\'on)$\times$(total de su columna).
\end{enumerate}
\end{receta}

\begin{ejemplo}{Ejemplo resuelto 1 --- L\'ineas de producci\'on \origen{Apuntes, Ej.~3.3}}
$X$ = art\'iculos de la l\'inea uno, $Y$ = art\'iculos de la l\'inea dos.
\begin{center}
\begin{tabular}{cc|ccc|c}
 & & \multicolumn{3}{c|}{$x$} & \\
 & $\fxy$ & 0 & 1 & 2 & $f_Y(y)$\\ \hline
 \multirow{3}{*}{$y$} & 0 & 0.10 & 0.20 & 0.20 & \textbf{0.5}\\
 & 1 & 0.04 & 0.08 & 0.08 & \textbf{0.2}\\
 & 2 & 0.06 & 0.12 & 0.12 & \textbf{0.3}\\ \hline
 & $f_X(x)$ & \textbf{0.20} & \textbf{0.40} & \textbf{0.40} & 1
\end{tabular}
\end{center}
\textbf{a) $P(X>Y)$.} Celdas con $x>y$: $(1,0)$, $(2,0)$, $(2,1)$:
\[P(X>Y)=0.20+0.20+0.08=\resp{0.48}\]
\textbf{b) Marginales.} Son los m\'argenes de la tabla (en negritas).

\textbf{c) Condicional de $X$ dado $Y=1$.} Rengl\'on $y=1$ entre su total $0.2$:
\begin{center}
\begin{tabular}{c|ccc}
$x$ & 0 & 1 & 2\\ \hline
$f_{X\mid Y=1}(x)$ & $0.04/0.2=0.2$ & $0.08/0.2=0.4$ & $0.08/0.2=0.4$
\end{tabular}
\end{center}
\textbf{d) $P(X<2\mid Y=1)$} $=0.2+0.4=\resp{0.6}$.

\textbf{e) ?`Independientes?} Cada celda es igual al producto de sus m\'argenes
($0.10=0.2\cdot0.5$, $0.08=0.4\cdot0.2$, \ldots, las nueve). \resp{\text{S\'i son independientes}}.
Observa que por eso la condicional del inciso c) sali\'o igual a la marginal $f_X$.
\end{ejemplo}

\begin{ejemplo}{Ejemplo resuelto 2 --- Supermercado \origen{Serie 3, problema 1}}
$X_1$ = clientes en caja r\'apida (renglones, 0 a 4), $X_2$ = clientes en caja superveloz (columnas, 0 a 3).
\begin{center}\small
\begin{tabular}{c|cccc|c}
$x_1\backslash x_2$ & 0 & 1 & 2 & 3 & $f_{X_1}$\\ \hline
0 & 0.08 & 0.07 & 0.04 & 0.00 & 0.19\\
1 & 0.06 & 0.15 & 0.05 & 0.04 & 0.30\\
2 & 0.05 & 0.04 & 0.10 & 0.06 & 0.25\\
3 & 0.00 & 0.03 & 0.04 & 0.07 & 0.14\\
4 & 0.00 & 0.01 & 0.05 & 0.06 & 0.12\\ \hline
$f_{X_2}$ & 0.19 & 0.30 & 0.28 & 0.23 & 1
\end{tabular}
\end{center}
\textbf{a)} $P(X_1=1,X_2=1)=\resp{0.15}$ (una sola celda).

\textbf{b)} $P(X_1=X_2)$ es la diagonal: $0.08+0.15+0.10+0.07=\resp{0.40}$.

\textbf{c)} $A=\{|X_1-X_2|\ge 2\}$. Celdas: $(0,2),(0,3),(1,3),(2,0),(3,0),(3,1),(4,0),(4,1),(4,2)$:
\[P(A)=0.04+0+0.04+0.05+0+0.03+0+0.01+0.05=\resp{0.22}\]
\textbf{d)} $P(X_1+X_2=4)$: celdas $(1,3),(2,2),(3,1),(4,0)$: $0.04+0.10+0.03+0=\resp{0.17}$.

$P(X_1+X_2\ge4)$: se agregan $(2,3),(3,2),(3,3),(4,1),(4,2),(4,3)$:
$0.17+0.06+0.04+0.07+0.01+0.05+0.06=\resp{0.46}$.
\end{ejemplo}

\begin{trampa}
\begin{itemize}
\item Leer la tabla al rev\'es. En la Serie~1 los renglones son $x_1$; en el Ej.~3.3 los renglones son $y$. Revisa siempre el encabezado.
\item \guillemotleft A lo sumo\guillemotright{} es $\le$, \guillemotleft por lo menos\guillemotright{} es $\ge$, \guillemotleft m\'as de\guillemotright{} es $>$.
\item Para \guillemotleft independientes\guillemotright{} basta \textbf{una} celda que falle para decir que no; para decir que s\'i debes revisar \textbf{todas}.
\end{itemize}
\end{trampa}

\begin{ejemplo}{Ejemplo resuelto 3 --- Restaurante \origen{Serie 3, problema 8}}
$X$ = costo del hombre (renglones), $Y$ = costo de la mujer (columnas).
\begin{center}
\begin{tabular}{c|ccc|c}
$x\backslash y$ & 7 & 9 & 10 & $f_X$\\ \hline
7 & 0.05 & 0.05 & 0.10 & 0.20\\
9 & 0.05 & 0.10 & 0.35 & 0.50\\
10 & 0 & 0.20 & 0.10 & 0.30\\ \hline
$f_Y$ & 0.10 & 0.35 & 0.55 & 1
\end{tabular}
\end{center}
\textbf{a)} Marginales: los m\'argenes.

\textbf{b)} $P(X\le9,\,Y\le9)=0.05+0.05+0.05+0.10=\resp{0.25}$.

\textbf{c)} $\fxy(10,7)=0$ pero $f_X(10)f_Y(7)=0.3\cdot0.1=0.03\neq0$. \resp{\text{No son independientes}}.

\textbf{d)} Costo total esperado: $\E(X+Y)=\E(X)+\E(Y)$.
\begin{align*}
\E(X)&=7(0.2)+9(0.5)+10(0.3)=8.90\\
\E(Y)&=7(0.10)+9(0.35)+10(0.55)=9.35
\end{align*}
$\E(X+Y)=\resp{18.25}$.

\textbf{e)} Devoluci\'on $=|X-Y|$. Se multiplica $|x-y|$ por cada celda (la diagonal aporta 0):
\[
\E|X-Y|=2(0.05)+3(0.10)+2(0.05)+1(0.35)+3(0)+1(0.20)=\resp{1.05}
\]
\end{ejemplo}

\begin{mecaniza}
Con la tabla del restaurante obt\'en la funci\'on condicional de $X$ dado $Y=9$ y $P(X\ge 9\mid Y=9)$.

\textbf{Respuesta:} columna $y=9$ entre $0.35$: $f_{X\mid Y=9}=\{7\!:\tfrac17,\;9\!:\tfrac27,\;10\!:\tfrac47\}$;
$P(X\ge9\mid Y=9)=\tfrac67\approx0.857$.
\end{mecaniza}

\subsection{Funci\'on de distribuci\'on conjunta (tabla acumulada)}

$F_{XY}(x,y)=P(X\le x,\,Y\le y)$. En una tabla: cada celda de $F$ es la suma de todas las celdas de $f$ que quedan
\textbf{arriba y a la izquierda} de ella, incluy\'endola.

\begin{ejemplo}{Ejemplo resuelto 4 \origen{Apuntes, Ej.~3.7 (en blanco en tus apuntes)}}
Con la tabla del Ej.~3.3:
\begin{center}
\begin{tabular}{cc|ccc}
 & $F_{XY}$ & $x=0$ & $x=1$ & $x=2$\\ \hline
 & $y=0$ & 0.10 & 0.30 & 0.50\\
 & $y=1$ & 0.14 & 0.42 & 0.70\\
 & $y=2$ & 0.20 & 0.60 & 1.00
\end{tabular}
\end{center}
Por ejemplo $F(1,1)=0.10+0.20+0.04+0.08=0.42$. Comprobaciones: la esquina final siempre vale 1; el \'ultimo
rengl\'on es $F_X$ y la \'ultima columna es $F_Y$.
\end{ejemplo}

\subsection{Construir la tabla desde el experimento}

\begin{ejemplo}{Ejemplo resuelto 5 --- Transistores \origen{Serie 3, problema 10a}}
Cinco transistores, tres defectuosos (D) y dos buenos (B). $X$ = pruebas hasta el primer D;
$Y$ = pruebas adicionales hasta el segundo D. Se multiplica rama por rama, sin reemplazo:
{\small\begin{align*}
P(X=1,Y=1)&=P(DD)=\tfrac35\cdot\tfrac24=\tfrac{3}{10} &
P(X=1,Y=2)&=P(DBD)=\tfrac35\cdot\tfrac24\cdot\tfrac23=\tfrac{2}{10}\\
P(X=1,Y=3)&=P(DBBD)=\tfrac35\cdot\tfrac24\cdot\tfrac13\cdot1=\tfrac{1}{10} &
P(X=2,Y=1)&=P(BDD)=\tfrac25\cdot\tfrac34\cdot\tfrac23=\tfrac{2}{10}\\
P(X=2,Y=2)&=P(BDBD)=\tfrac25\cdot\tfrac34\cdot\tfrac13\cdot1=\tfrac{1}{10} &
P(X=3,Y=1)&=P(BBDD)=\tfrac25\cdot\tfrac14\cdot1\cdot1=\tfrac{1}{10}
\end{align*}}
\begin{center}
\begin{tabular}{c|ccc|c}
$x\backslash y$ & 1 & 2 & 3 & $f_X$\\ \hline
1 & 0.3 & 0.2 & 0.1 & 0.6\\
2 & 0.2 & 0.1 & 0 & 0.3\\
3 & 0.1 & 0 & 0 & 0.1\\ \hline
$f_Y$ & 0.6 & 0.3 & 0.1 & 1
\end{tabular}
\end{center}
La suma da 1, as\'i que no falta ning\'un caso. Esta tabla se vuelve a usar en la secci\'on de correlaci\'on.
\end{ejemplo}

\section{Variables conjuntas continuas (la densidad)}

Ahora $\fxy(x,y)$ es una superficie y la probabilidad es el \textbf{volumen} bajo ella sobre una regi\'on.
Todo el trabajo est\'a en una sola habilidad: \textbf{poner bien los l\'imites de la integral doble}.

\subsection{La habilidad clave: l\'imites de integraci\'on}

\begin{receta}{poner l\'imites (m\'etodo de la flecha)}
\begin{enumerate}
\item \textbf{Dibuja la regi\'on} donde la densidad no es cero (el \guillemotleft soporte\guillemotright).
Encima, sombrea la parte que pide el evento.
\item Elige la variable de \textbf{afuera}. Sus l\'imites son \textbf{n\'umeros}: el valor m\'inimo y m\'aximo que toma en la zona sombreada.
\item Para la variable de \textbf{adentro}, traza una flecha que cruce la zona en su direcci\'on
(vertical si adentro es $y$, horizontal si adentro es $x$). Donde la flecha \textbf{entra} es el l\'imite inferior;
donde \textbf{sale}, el superior. Pueden depender de la variable de afuera.
\item Si la flecha entra o sale por curvas distintas seg\'un d\'onde la pongas, parte la integral en dos.
\item Integra primero adentro (la otra variable es constante) y luego afuera.
\end{enumerate}
\end{receta}

\begin{center}
\begin{tikzpicture}[scale=2.35]
% rectangulo
\begin{scope}
\draw[->] (-0.1,0)--(1.25,0) node[right]{$x$}; \draw[->] (0,-0.1)--(0,1.25) node[above]{$y$};
\fill[azul!18] (0,0) rectangle (1,1); \draw (0,0) rectangle (1,1);
\draw[->,rojo,thick] (0.55,-0.05)--(0.55,1.08);
\node[below] at (1,0){\small 1}; \node[left] at (0,1){\small 1};
\node[align=center] at (0.5,-0.5){\small Rect\'angulo\\[3pt]\small $\displaystyle\int_0^1\!\!\int_0^1$};
\end{scope}
% triangulo x+y<=1
\begin{scope}[xshift=1.85cm]
\draw[->] (-0.1,0)--(1.25,0) node[right]{$x$}; \draw[->] (0,-0.1)--(0,1.25) node[above]{$y$};
\fill[azul!18] (0,0)--(1,0)--(0,1)--cycle; \draw (0,0)--(1,0)--(0,1)--cycle;
\draw[->,rojo,thick] (0.35,-0.05)--(0.35,0.73);
\node[right] at (0.5,0.55){\small $y=1-x$};
\node[align=center] at (0.5,-0.5){\small $x+y\le1$\\[3pt]\small $\displaystyle\int_0^1\!\!\int_0^{1-x}dy\,dx$};
\end{scope}
% triangulo y<=x
\begin{scope}[xshift=3.7cm]
\draw[->] (-0.1,0)--(1.25,0) node[right]{$x$}; \draw[->] (0,-0.1)--(0,1.25) node[above]{$y$};
\fill[azul!18] (0,0)--(1,0)--(1,1)--cycle; \draw (0,0)--(1,0)--(1,1)--cycle;
\draw[->,rojo,thick] (0.65,-0.05)--(0.65,0.73);
\node[left] at (0.5,0.6){\small $y=x$};
\node[align=center] at (0.5,-0.5){\small $0\le y\le x\le1$\\[3pt]\small $\displaystyle\int_0^1\!\!\int_0^{x}dy\,dx$};
\end{scope}
% triangulo x<=y
\begin{scope}[xshift=5.55cm]
\draw[->] (-0.1,0)--(1.25,0) node[right]{$x$}; \draw[->] (0,-0.1)--(0,1.25) node[above]{$y$};
\fill[azul!18] (0,0)--(0,1)--(1,1)--cycle; \draw (0,0)--(0,1)--(1,1)--cycle;
\draw[->,rojo,thick] (0.35,0.3)--(0.35,1.08);
\node[right] at (0.55,0.5){\small $y=x$};
\node[align=center] at (0.5,-0.5){\small $0\le x\le y\le1$\\[3pt]\small $\displaystyle\int_0^1\!\!\int_x^{1}dy\,dx$};
\end{scope}
\end{tikzpicture}
\end{center}
Estas cuatro regiones cubren pr\'acticamente todos los problemas de tus apuntes y de la serie. Apr\'endetelas de vista.

\subsection{Hallar la constante y calcular probabilidades}

\begin{receta}{hallar $k$}
Plantea $\displaystyle\iint_{\text{soporte}} \fxy\,dx\,dy=1$, integra y despeja $k$. Despu\'es \textbf{reescribe} la densidad ya con el valor de $k$.
\end{receta}

\begin{ejemplo}{Ejemplo resuelto 6 \origen{Apuntes, Ej.~3.4}}
$\fxy(x,y)=k(x+y)$ para $0\le x\le1$, $0\le y\le1$.

\textbf{a) $k$.}
\[
1=k\int_0^1\!\!\int_0^1 (x+y)\,dx\,dy=k\int_0^1\Big(\tfrac12+y\Big)dy=k\Big(\tfrac12+\tfrac12\Big)\;\Rightarrow\;\resp{k=1}
\]
\textbf{b) $P(0.5\le X\le1,\;0.5\le Y\le1)$.} Regi\'on rectangular, l\'imites num\'ericos:
\[
\int_{0.5}^1\!\!\int_{0.5}^1 (x+y)\,dx\,dy=\int_{0.5}^1\Big(\tfrac{y}{2}+\tfrac38\Big)dy=\Big[\tfrac{y^2}{4}+\tfrac38y\Big]_{0.5}^{1}=\resp{0.375}
\]
\textbf{c) $P(X+Y\le1)$.} Es el tri\'angulo bajo la recta $y=1-x$ (segunda figura de arriba):
\[
\int_0^1\!\!\int_0^{1-x}(x+y)\,dy\,dx=\int_0^1\Big[x(1-x)+\tfrac{(1-x)^2}{2}\Big]dx=\int_0^1\frac{1-x^2}{2}\,dx=\tfrac12-\tfrac16=\resp{\tfrac13}
\]
\end{ejemplo}

\begin{ejemplo}{Ejemplo resuelto 7 --- Tabacalera \origen{Serie 3, problema 3a}}
$\fxy(x,y)=24xy$ para $0\le x\le1$, $0\le y\le1$, $x+y\le1$. Calcular $P(X>\tfrac12)$.

\begin{minipage}{0.29\textwidth}
\begin{tikzpicture}[scale=2.4]
\draw[->] (-0.1,0)--(1.25,0) node[right]{$x$}; \draw[->] (0,-0.1)--(0,1.25) node[above]{$y$};
\draw (0,0)--(1,0)--(0,1)--cycle;
\fill[azul!25] (0.5,0)--(1,0)--(0.5,0.5)--cycle;
\draw[dashed] (0.5,0)--(0.5,0.5);
\draw[->,rojo,thick] (0.7,-0.05)--(0.7,0.38);
\node[below] at (0.5,0){\small $\tfrac12$}; \node[below] at (1,0){\small 1};
\node[right] at (0.45,0.65){\small $y=1-x$};
\end{tikzpicture}
\end{minipage}
\begin{minipage}{0.69\textwidth}\small
Afuera $x$: de $\tfrac12$ a $1$. La flecha vertical entra en $y=0$ y sale en $y=1-x$.
\begin{align*}
P\big(X>\tfrac12\big)&=\int_{1/2}^{1}\!\int_0^{1-x}24xy\,dy\,dx=\int_{1/2}^{1}12x(1-x)^2dx\\
&=12\Big[\tfrac{x^2}{2}-\tfrac{2x^3}{3}+\tfrac{x^4}{4}\Big]_{1/2}^{1}
=12\Big(\tfrac{16}{192}-\tfrac{11}{192}\Big)=\resp{\tfrac{5}{16}=0.3125}
\end{align*}
\end{minipage}
\end{ejemplo}

\begin{ejemplo}{Ejemplo resuelto 8 --- Regi\'on en banda, integral partida \origen{Serie 3, problema 2}}
$\fxy(x,y)=kx$ para $0\le x\le4$, $y\ge0$, $-x+2\le y\le -x+4$.

\begin{minipage}{0.36\textwidth}
\begin{tikzpicture}[scale=0.82]
\draw[->] (-0.2,0)--(4.7,0) node[right]{$x$}; \draw[->] (0,-0.2)--(0,4.6) node[above]{$y$};
\fill[azul!18] (0,2)--(0,4)--(4,0)--(2,0)--cycle;
\draw (0,2)--(0,4)--(4,0)--(2,0)--cycle;
\draw[dashed] (2,0)--(2,2);
\draw[->,rojo,thick] (1,0.85)--(1,3.2);
\draw[->,rojo,thick] (3,-0.15)--(3,1.2);
\node[below] at (2,0){\small 2}; \node[below] at (4,0){\small 4};
\node[left] at (0,2){\small 2}; \node[left] at (0,4){\small 4};
\node[right] at (2.1,2.3){\small $y=4-x$};
\node[left] at (1.25,0.45){\small $y=2-x$};
\end{tikzpicture}
\end{minipage}
\begin{minipage}{0.62\textwidth}
La flecha vertical entra por curvas distintas:
\begin{itemize}
\item si $0\le x\le2$: entra en $y=2-x$, sale en $y=4-x$;
\item si $2\le x\le4$: entra en $y=0$, sale en $y=4-x$.
\end{itemize}
Por eso la integral se parte en $x=2$.
\end{minipage}

\textbf{a) $k$.}
\[
1=k\Big[\int_0^2\!\!\int_{2-x}^{4-x}x\,dy\,dx+\int_2^4\!\!\int_0^{4-x}x\,dy\,dx\Big]
=k\Big[\int_0^2 2x\,dx+\int_2^4(4x-x^2)dx\Big]=k\Big[4+\tfrac{16}{3}\Big]
\]
$\Rightarrow\ \resp{k=\tfrac{3}{28}}$.

\textbf{b) $P(1\le X\le3)$.} Mismas dos piezas, ahora con $x$ de 1 a 2 y de 2 a 3:
\[
\tfrac{3}{28}\Big[\int_1^2 2x\,dx+\int_2^3(4x-x^2)dx\Big]=\tfrac{3}{28}\Big[3+\tfrac{11}{3}\Big]=\tfrac{3}{28}\cdot\tfrac{20}{3}=\resp{\tfrac57\approx0.714}
\]
\end{ejemplo}

\begin{ejemplo}{Ejemplo resuelto 9 --- Densidad constante: usa \'areas \origen{Apuntes, Ej.~5.5}}
$f(x_1,x_2)=2$ en el tri\'angulo $x_1\ge0$, $x_2\ge0$, $x_1+x_2\le1$.
Si la densidad es constante, \textbf{probabilidad $=$ constante $\times$ \'area}; no hace falta integrar.

\textbf{a)} $P(X_1\le\tfrac34,X_2\le\tfrac34)$: cuadrado de lado $\tfrac34$ menos la esquina que queda arriba de la recta
(tri\'angulo de catetos $\tfrac12$): \'area $=\tfrac{9}{16}-\tfrac18=\tfrac{7}{16}$. Probabilidad $=2\cdot\tfrac{7}{16}=\resp{\tfrac78}$.

\textbf{b)} $P(X_1\le\tfrac12,X_2\le\tfrac12)$: cuadrado de lado $\tfrac12$ completo dentro del tri\'angulo: $2\cdot\tfrac14=\resp{\tfrac12}$.

\textbf{c)} $P(X_1\le\tfrac12\mid X_2\le\tfrac12)=\dfrac{P(X_1\le\frac12,\,X_2\le\frac12)}{P(X_2\le\frac12)}$.
El denominador es $\int_0^{1/2}\!\int_0^{1-x_2}2\,dx_1dx_2=\tfrac34$. Resultado: $\dfrac{1/2}{3/4}=\resp{\tfrac23}$.

\textbf{Ojo:} aqu\'i se condiciona a un \emph{evento} ($X_2\le\frac12$), as\'i que se usa $P(A\cap B)/P(B)$,
no la densidad condicional.
\end{ejemplo}

\begin{ejemplo}{Ejemplo resuelto 10 --- Neum\'aticos \origen{Serie 3, problema 7}}
$\fxy(x,y)=K(x^2+y^2)$ para $20\le x\le30$, $20\le y\le30$.

\textbf{a)} Por simetr\'ia $\iint(x^2+y^2)=2\cdot10\int_{20}^{30}x^2dx=20\cdot\tfrac{19000}{3}=\tfrac{380000}{3}$, as\'i que
$\resp{K=\tfrac{3}{380000}}$.

\textbf{b)} Ambos menores que 26:
\[
K\int_{20}^{26}\!\!\int_{20}^{26}(x^2+y^2)\,dx\,dy=K\cdot2\cdot6\cdot\frac{26^3-20^3}{3}=\frac{3}{380000}(38304)=\resp{0.3024}
\]
\textbf{c)} $P(|X-Y|\le2)$. Es la banda alrededor de la diagonal. Conviene el complemento: las dos esquinas son
iguales por simetr\'ia.

\begin{minipage}{0.3\textwidth}
\begin{tikzpicture}[scale=0.36]
\draw[->] (19,20)--(31.5,20) node[right]{\small $x$}; \draw[->] (20,19)--(20,31.5) node[above]{\small $y$};
\draw (20,20) rectangle (30,30);
\fill[azul!22] (20,20)--(22,20)--(30,28)--(30,30)--(28,30)--(20,22)--cycle;
\fill[rojo!20] (20,22)--(28,30)--(20,30)--cycle;
\fill[rojo!20] (22,20)--(30,28)--(30,20)--cycle;
\node at (27.2,22.4){\small $y<x-2$};
\node at (22.8,27.6){\small $y>x+2$};
\node[below] at (20,20){\small 20}; \node[below] at (30,20){\small 30};
\end{tikzpicture}
\end{minipage}
\begin{minipage}{0.68\textwidth}
\[
P(|X-Y|\le2)=1-2K\int_{20}^{28}\!\!\int_{x+2}^{30}(x^2+y^2)\,dy\,dx=\resp{0.3593}
\]
(En el examen el planteamiento con la figura vale casi todos los puntos.)
\end{minipage}

\textbf{d)} Marginal del neum\'atico derecho:
\[
f_X(x)=K\int_{20}^{30}(x^2+y^2)\,dy=K\Big(10x^2+\tfrac{19000}{3}\Big)=\resp{\frac{3x^2}{38000}+\frac1{20}},\quad 20\le x\le30
\]
\textbf{e)} $f_X(x)f_Y(y)$ tiene un t\'ermino $x^2y^2$ que la conjunta no tiene: $\fxy\neq f_Xf_Y$. \resp{\text{No son independientes}}.
\end{ejemplo}

\begin{mecaniza}
\textbf{1.} (Autoexamen 1) $\fxy=c$ en $0\le x\le0.25$, $0\le y\le2000$. Halla $c$.
\quad\textbf{R:} $c\cdot0.25\cdot2000=1\Rightarrow c=\tfrac1{500}$ (inciso B).

\textbf{2.} (Autoexamen 2) $\fxy=6x^2y$ en el cuadrado unitario. $P(X\le0.5,Y\le0.5)$.
\quad\textbf{R:} $6\cdot\tfrac{(0.5)^3}{3}\cdot\tfrac{(0.5)^2}{2}=\tfrac1{32}=0.03125$ (inciso A).

\textbf{3.} Con $\fxy=24xy$ del tri\'angulo, calcula $P(X\le\tfrac12)$.\quad\textbf{R:} $1-\tfrac{5}{16}=\tfrac{11}{16}$.
\end{mecaniza}

\subsection{Marginales y condicionales continuas}

\begin{receta}{marginal y condicional de una densidad}
\begin{enumerate}
\item $f_X(x)=\int \fxy\,dy$: integra en $y$ con los l\'imites que da la flecha vertical. \textbf{Escribe el intervalo de $x$} donde vale.
\item $f_Y(y)=\int \fxy\,dx$: flecha horizontal. Escribe el intervalo de $y$.
\item Condicional: $f_{X\mid y_0}(x\mid y_0)=\dfrac{\fxy(x,y_0)}{f_Y(y_0)}$, con el intervalo de $x$ que queda \textbf{cuando $y=y_0$}.
\item Probabilidad condicional con valor exacto ($P(X<a\mid Y=y_0)$): sustituye $y_0$ en la condicional e integra en $x$.
\end{enumerate}
\end{receta}

\begin{ejemplo}{Ejemplo resuelto 11 \origen{Apuntes, Ej.~3.6 (en blanco) y Ej.~3.10}}
$\fxy(x,y)=x^3+\dfrac{xy^2}{a}$ para $0\le x\le1$, $0\le y\le2$.

\textbf{a) Valor de $a$.}
\[
1=\int_0^2\!\!\int_0^1\Big(x^3+\frac{xy^2}{a}\Big)dx\,dy=\int_0^2\Big(\frac14+\frac{y^2}{2a}\Big)dy=\frac12+\frac{4}{3a}
\;\Rightarrow\;\resp{a=\tfrac83}
\]
es decir $\fxy=x^3+\tfrac38xy^2$ (la misma del Ej.~3.10).

\textbf{b) $P(X\le Y)$.} Flecha vertical: entra en $y=x$, sale en $y=2$; $x$ de 0 a 1.
\[
\int_0^1\!\!\int_x^2\Big(x^3+\tfrac38xy^2\Big)dy\,dx=\int_0^1\Big(x+2x^3-\tfrac98x^4\Big)dx=\tfrac12+\tfrac12-\tfrac{9}{40}=\resp{\tfrac{31}{40}=0.775}
\]
\textbf{c) Marginales.}
\begin{align*}
f_X(x)&=\int_0^2\Big(x^3+\tfrac38xy^2\Big)dy=\resp{2x^3+x},\quad 0\le x\le1\\
f_Y(y)&=\int_0^1\Big(x^3+\tfrac38xy^2\Big)dx=\resp{\tfrac14+\tfrac{3}{16}y^2},\quad 0\le y\le2
\end{align*}
\textbf{d) Condicionales.}
\begin{align*}
f_{X\mid y_0}(x\mid y_0)&=\frac{x^3+\frac38xy_0^2}{\frac14+\frac{3}{16}y_0^2}=\resp{\frac{2x(8x^2+3y_0^2)}{4+3y_0^2}},\quad 0\le x\le1\\[3pt]
f_{Y\mid x_0}(y\mid x_0)&=\frac{x_0^3+\frac38x_0y^2}{2x_0^3+x_0}=\resp{\frac{8x_0^2+3y^2}{8(2x_0^2+1)}},\quad 0\le y\le2
\end{align*}
(coinciden con las que da el Ej.~3.10 de tus apuntes).
\end{ejemplo}

\begin{ejemplo}{Ejemplo resuelto 12 --- Tabacalera, incisos b y c \origen{Serie 3, problema 3}}
$\fxy=24xy$ en el tri\'angulo $x+y\le1$.

\textbf{b) Marginal de $Y$.} Flecha horizontal: entra en $x=0$, sale en $x=1-y$.
\[
f_Y(y)=\int_0^{1-y}24xy\,dx=\resp{12y(1-y)^2},\quad0\le y\le1
\]
\textbf{c) $P(X<\tfrac18\mid Y=\tfrac34)$.} Condicional con $y_0=\tfrac34$; ah\'i $x$ s\'olo puede ir de 0 a $1-\tfrac34=\tfrac14$:
\[
f_{X\mid 3/4}(x)=\frac{24x\cdot\frac34}{12\cdot\frac34\cdot\big(\frac14\big)^2}=32x,\quad0\le x\le\tfrac14
\qquad\Rightarrow\qquad
P=\int_0^{1/8}32x\,dx=16\cdot\tfrac{1}{64}=\resp{0.25}
\]
\end{ejemplo}

\begin{trampa}
\begin{itemize}
\item En regiones triangulares los l\'imites de la marginal \textbf{no son} $0$ y $1$: son los de la flecha ($0$ a $1-y$, $y$ a $1$, etc.).
\item La marginal de $X$ no puede contener $y$. Si te queda una $y$, los l\'imites est\'an mal.
\item Al condicionar en un tri\'angulo, el intervalo de la variable libre cambia con el valor dado (arriba: $0\le x\le\frac14$, no $0\le x\le1$).
\item Comprueba: toda marginal y toda condicional debe integrar 1.
\end{itemize}
\end{trampa}

\begin{mecaniza}
\textbf{1.} (Autoexamen 3) $\fxy=\dfrac{3x-y}{18}$, $1<x<4$, $1<y<2$. Marginal de $X$.
\quad\textbf{R:} $\int_1^2\frac{3x-y}{18}dy=\frac{x}{6}-\frac{1}{12}$ (inciso B).

\textbf{2.} (Autoexamen 4) $\fxy=8xy$, $0<x<y<1$. Condicional de $X$ dado $y$.
\quad\textbf{R:} $f_Y(y)=\int_0^y8xy\,dx=4y^3$; $f_{X\mid y}=\dfrac{8xy}{4y^3}=\dfrac{2x}{y^2}$, $0<x<y$ (inciso B).

\textbf{3.} (Serie 6) $\fxy=\tfrac1{240}$ en $8.5\le x\le10.5$, $120\le y\le240$. Obt\'en $f_{X\mid y}$ y $f_{Y\mid x}$.
\quad\textbf{R:} $f_X=\tfrac12$, $f_Y=\tfrac1{120}$; $f_{X\mid y}=\tfrac{1/240}{1/120}=\tfrac12=f_X$ y
$f_{Y\mid x}=\tfrac1{120}=f_Y$. Son iguales a las marginales porque hay independencia: saber el colesterol no dice nada del calcio.
\end{mecaniza}

\subsection{Independencia en el caso continuo}

\begin{receta}{decidir independencia r\'apido}
\begin{enumerate}
\item \textbf{Mira el soporte.} Si no es un rect\'angulo (los l\'imites de una variable dependen de la otra:
tri\'angulos, $y\le x$, $x+y\le1$), son \textbf{dependientes}. Justificaci\'on: hay puntos donde $\fxy=0$ pero $f_Xf_Y>0$.
\item Si el soporte es rect\'angulo, calcula $f_X$ y $f_Y$ y compara $f_Xf_Y$ con $\fxy$.
Atajo: si $\fxy$ se puede escribir como (algo s\'olo de $x$)$\times$(algo s\'olo de $y$), son independientes; si es una suma como $x+y$ o $x^2+y^2$, no.
\item En el examen \textbf{escribe la comparaci\'on}: calcula las marginales y muestra el producto.
\end{enumerate}
\end{receta}

\begin{ejemplo}{Ejemplo resuelto 13 \origen{Apuntes, Ej.~3.8 (en blanco)}}
$\fxy=\tfrac83x(x+y)$ para $0\le x\le1$, $0\le y\le x$.
\begin{align*}
f_X(x)&=\int_0^x\tfrac83(x^2+xy)\,dy=\tfrac83\Big(x^3+\tfrac{x^3}{2}\Big)=4x^3,\quad0\le x\le1\\
f_Y(y)&=\int_y^1\tfrac83(x^2+xy)\,dx=\tfrac89(1-y^3)+\tfrac43y(1-y^2)=\tfrac89+\tfrac43y-\tfrac{20}{9}y^3,\quad0\le y\le1
\end{align*}
$f_X(x)f_Y(y)=4x^3\big(\tfrac89+\tfrac43y-\tfrac{20}{9}y^3\big)\neq\tfrac83x(x+y)$. \resp{\text{No son independientes}}
(adem\'as el soporte es triangular).
\end{ejemplo}

\begin{ejemplo}{Ejemplo resuelto 14 \origen{Apuntes, Ej.~3.12 (en blanco), incisos a y b}}
$\fxy=4xye^{-(x^2+y^2)}$ para $x\ge0$, $y\ge0$.

Con $u=y^2$: $\int_0^\infty 2ye^{-y^2}dy=1$. Entonces
\[
f_X(x)=2xe^{-x^2}\int_0^\infty2ye^{-y^2}dy=\resp{2xe^{-x^2}},\ x\ge0\qquad f_Y(y)=\resp{2ye^{-y^2}},\ y\ge0
\]
$f_Xf_Y=4xye^{-(x^2+y^2)}=\fxy$: \resp{\text{independientes}}. Por lo tanto las condicionales son las marginales:
$f_{X\mid y}=2xe^{-x^2}$ y $f_{Y\mid x}=2ye^{-y^2}$. (Los incisos c y d se terminan en la secci\'on~\ref{sec:cov}.)
\end{ejemplo}

\section{Valor esperado de funciones de $X$ y $Y$}

\begin{teoria}{Definici\'on y propiedades}
\[
\E[g(X,Y)]=\sum_x\sum_y g(x,y)\,\fxy(x,y)\qquad\text{o}\qquad\E[g(X,Y)]=\iint g(x,y)\,\fxy(x,y)\,dx\,dy
\]
\guillemotleft Multiplica lo que quieres promediar por la conjunta y suma (o integra) sobre todo el soporte.\guillemotright
\begin{itemize}
\item $\E(c)=c$ \qquad $\E(aX+bY+c)=a\E(X)+b\E(Y)+c$ \quad(\textbf{siempre}, sean o no independientes).
\item S\'olo si son independientes: $\E(XY)=\E(X)\,\E(Y)$.
\end{itemize}
Casos que vas a usar: $g=x$ da $\E(X)$; $g=y$ da $\E(Y)$; $g=xy$ da $\E(XY)$; $g=x^2$ da $\E(X^2)$.
\end{teoria}

\begin{receta}{traducir el enunciado a $g(X,Y)$}
\begin{center}
\begin{tabular}{ll}
\guillemotleft total\guillemotright, \guillemotleft entre los dos\guillemotright & $g=X+Y$\\
\guillemotleft lo que queda\guillemotright, \guillemotleft sobrante\guillemotright & $g=Y-X$ (lo que hab\'ia menos lo que sali\'o)\\
\guillemotleft diferencia entre el m\'as caro y el m\'as barato\guillemotright & $g=|X-Y|$\\
\guillemotleft proporci\'on de una proporci\'on\guillemotright & $g=XY$\\
\guillemotleft costo $aX+b$\guillemotright & $g=aX+b$
\end{tabular}
\end{center}
\end{receta}

\begin{ejemplo}{Ejemplo resuelto 15 --- Queroseno \origen{Serie 3, problema 5}}
$Y$ = cantidad al inicio, $X$ = cantidad vendida, $\fxy=2$ en $0\le x\le y\le1$. Lo que queda es $Y-X$.
\[
\E(Y-X)=\int_0^1\!\!\int_0^y2(y-x)\,dx\,dy=\int_0^1\big[2yx-x^2\big]_0^{y}dy=\int_0^1y^2\,dy=\resp{\tfrac13}
\]
Quedan en promedio $\tfrac13$ de mil litros, unos 333 litros.
\end{ejemplo}

\begin{ejemplo}{Ejemplo resuelto 16 --- Impurezas \origen{Apuntes, Ej.~3.9}}
$\fxy=2-2x$ en el cuadrado unitario. La proporci\'on de impureza tipo I en la muestra es $XY$:
\[
\E(XY)=\int_0^1\!\!\int_0^1xy(2-2x)\,dy\,dx=\int_0^1x(1-x)\,dx=\tfrac12-\tfrac13=\resp{\tfrac16}
\]
\end{ejemplo}

\begin{ejemplo}{Ejemplo resuelto 17 --- Tabla independiente \origen{Serie 3, problema 4}}
\begin{center}
\begin{tabular}{c|cc|c}
$y\backslash x$ & 2 & 4 & $f_Y$\\ \hline
1 & 0.10 & 0.15 & 0.25\\
3 & 0.20 & 0.30 & 0.50\\
5 & 0.10 & 0.15 & 0.25\\ \hline
$f_X$ & 0.40 & 0.60 & 1
\end{tabular}
\end{center}
$\E(X)=2(0.4)+4(0.6)=3.2$,\quad $\E(Y)=1(0.25)+3(0.5)+5(0.25)=3$.

\textbf{a)} $\E(2X-3Y)=2(3.2)-3(3)=\resp{-2.6}$.

\textbf{b)} Como son independientes, $\E(XY)=\E(X)\E(Y)=3.2\cdot3=\resp{9.6}$.
(Comprobaci\'on directa: $2\cdot1\cdot0.10+4\cdot1\cdot0.15+2\cdot3\cdot0.20+4\cdot3\cdot0.30+2\cdot5\cdot0.10+4\cdot5\cdot0.15=9.6$.)
\end{ejemplo}

\begin{mecaniza}
\textbf{1.} (Autoexamen 5) Tabla: $f(0,1)=0.1$, $f(0,2)=0.3$, $f(1,1)=0.4$, $f(1,2)=0.2$. $\E(XY)$.
\quad\textbf{R:} $1\cdot1\cdot0.4+1\cdot2\cdot0.2=0.8$ (inciso B).

\textbf{2.} (Autoexamen 6) Dos l\'ineas de ensamble; n\'umero esperado de piezas de las dos.
$f_X=(0.35,0.37,0.28)$, $f_Y=(0.35,0.35,0.30)$.
\quad\textbf{R:} $\E(X)+\E(Y)=0.93+0.95=1.88$ (inciso D).
\end{mecaniza}

\subsection{Valor esperado condicional y curva de regresi\'on}

\begin{teoria}{Curva de regresi\'on}
Es el promedio de una variable cuando ya conozco la otra, usando la \textbf{condicional}:
\[
\hat y=\mu_{Y\mid x}=\E[Y\mid X=x]=\int y\,f_{Y\mid x}(y\mid x)\,dy
\qquad
\hat x=\mu_{X\mid y}=\E[X\mid Y=y]=\int x\,f_{X\mid y}(x\mid y)\,dx
\]
(con sumas si es discreta). El resultado es una funci\'on de la variable conocida; su gr\'afica es la curva de regresi\'on.
Sirve para pronosticar. Si son independientes, la curva es una constante: $\E[Y\mid x]=\E(Y)$.
\end{teoria}

\begin{receta}{curva de regresi\'on de $Y$ dado $x$}
(1) marginal $f_X$;\quad (2) condicional $f_{Y\mid x}=\fxy/f_X$;\quad (3) $\int y\,f_{Y\mid x}\,dy$ con los l\'imites de $y$;\quad
(4) simplifica y di si es una recta.
\end{receta}

\begin{ejemplo}{Ejemplo resuelto 18 \origen{Apuntes, Ej.~3.10 (en blanco)}}
Con $\fxy=x^3+\tfrac38xy^2$ y las condicionales del ejemplo resuelto 11:
\begin{align*}
\hat y=\E[Y\mid x]&=\int_0^2y\,\frac{8x^2+3y^2}{8(2x^2+1)}\,dy=\frac{16x^2+12}{8(2x^2+1)}=\resp{\frac{4x^2+3}{2(2x^2+1)}},\quad0\le x\le1\\[3pt]
\hat x=\E[X\mid y]&=\int_0^1x\,\frac{2x(8x^2+3y^2)}{4+3y^2}\,dx=\frac{\frac{16}{5}+2y^2}{4+3y^2}=\resp{\frac{2(5y^2+8)}{5(3y^2+4)}},\quad0\le y\le2
\end{align*}
No son rectas.
\end{ejemplo}

\begin{ejemplo}{Ejemplo resuelto 19 --- Regresi\'on lineal \origen{molde del Ej.~3.8}}
$\fxy=\tfrac83x(x+y)$, $0\le y\le x\le1$, con $f_X=4x^3$.
\[
f_{Y\mid x}=\frac{\frac83x(x+y)}{4x^3}=\frac{2(x+y)}{3x^2},\ 0\le y\le x
\qquad
\E[Y\mid x]=\int_0^xy\,\frac{2(x+y)}{3x^2}\,dy=\frac{2}{3x^2}\Big(\frac{x^3}{2}+\frac{x^3}{3}\Big)=\resp{\tfrac59x}
\]
Es una recta que pasa por el origen.
\end{ejemplo}

\begin{mecaniza}
\textbf{1.} (Serie 6c) Calcio y colesterol, densidad uniforme. Curvas de regresi\'on.
\quad\textbf{R:} hay independencia, as\'i que $\hat x=\E(X)=9.5$ y $\hat y=\E(Y)=180$: rectas constantes (una vertical y una horizontal).

\textbf{2.} (Serie 5) Con $\fxy=2$, $0\le x\le y\le1$: $\E[X\mid y]$ y $\E[Y\mid x]$.
\quad\textbf{R:} $f_Y=2y$, $f_{X\mid y}=\tfrac1y$ en $0\le x\le y$, $\E[X\mid y]=\tfrac y2$.
$f_X=2(1-x)$, $f_{Y\mid x}=\tfrac1{1-x}$ en $x\le y\le1$, $\E[Y\mid x]=\tfrac{1+x}2$. Ambas son rectas.

\textbf{3.} (Discreto, tabla del Ej.~3.3) $\E[X\mid Y=1]$.
\quad\textbf{R:} $0(0.2)+1(0.4)+2(0.4)=1.2$.
\end{mecaniza}

\subsection{Generadora de momentos de una suma (poco probable, pero es un rengl\'on)}

Si $X_1,\dots,X_n$ son \textbf{independientes} y $Y=X_1+\dots+X_n$, entonces
$M_Y(\theta)=M_{X_1}(\theta)\cdots M_{X_n}(\theta)$.

Ej.~3.11: $M_X=\dfrac{2}{2-\theta}$, $M_Y=\dfrac{1}{1-\theta}$, $T=X+Y$ $\Rightarrow$
$M_T(\theta)=\resp{\dfrac{2}{(2-\theta)(1-\theta)}}$.

\section{Covariancia y correlaci\'on}\label{sec:cov}

\begin{teoria}{Qu\'e miden}
\[
\Cov(X,Y)=\sigma_{XY}=\E(XY)-\E(X)\,\E(Y)
\qquad
\rho=\frac{\Cov(X,Y)}{\sqrt{\Var(X)}\sqrt{\Var(Y)}}=\frac{\sigma_{XY}}{\sigma_X\sigma_Y}
\qquad
\rho^2=\frac{\sigma_{XY}^2}{\sigma_X^2\sigma_Y^2}
\]
\begin{itemize}
\item \textbf{Signo de la covariancia:} positiva si cuando una sube la otra tiende a subir; negativa si una sube y la otra baja.
\item $\rho$ es la covariancia sin unidades: $-1\le\rho\le1$. Mide qu\'e tan \textbf{lineal} es la relaci\'on.
$\rho=\pm1$ s\'olo si $Y=mX+b$ exactamente.
\item $\rho^2$ (coeficiente de determinaci\'on): qu\'e fracci\'on de la variaci\'on de $Y$ se explica conociendo $X$.
\item Independientes $\Rightarrow\Cov=0$ y $\rho=0$. Al rev\'es no vale: $\Cov=0$ no garantiza independencia.
\end{itemize}
\end{teoria}

\begin{receta}{covariancia y correlaci\'on (las 5 esperanzas)}
\begin{enumerate}
\item Calcula $\E(X)$, $\E(Y)$, $\E(X^2)$, $\E(Y^2)$, $\E(XY)$.
(Las cuatro primeras salen m\'as r\'apido con las marginales.)
\item $\Var(X)=\E(X^2)-[\E(X)]^2$,\quad $\Var(Y)=\E(Y^2)-[\E(Y)]^2$.
\item $\Cov=\E(XY)-\E(X)\E(Y)$.
\item $\rho=\Cov/(\sigma_X\sigma_Y)$. Verifica que quede entre $-1$ y $1$.
\end{enumerate}
\end{receta}

\begin{ejemplo}{Ejemplo resuelto 20 --- Transistores \origen{Serie 3, problema 10b}}
Con la tabla del ejemplo resuelto 5: $f_X=f_Y=(0.6,\,0.3,\,0.1)$ en los valores $1,2,3$.
\begin{align*}
\E(X)&=\E(Y)=1(0.6)+2(0.3)+3(0.1)=1.5\\
\E(X^2)&=\E(Y^2)=1(0.6)+4(0.3)+9(0.1)=2.7\\
\Var(X)&=\Var(Y)=2.7-(1.5)^2=0.45\\
\E(XY)&=1(0.3)+2(0.2)+3(0.1)+2(0.2)+4(0.1)+3(0.1)=2.1\\
\Cov(X,Y)&=2.1-(1.5)(1.5)=-0.15\qquad\qquad\rho=\frac{-0.15}{\sqrt{0.45}\sqrt{0.45}}=\resp{-\tfrac13}
\end{align*}
Tiene sentido que sea negativa: si el primer defectuoso tarda en salir, quedan menos piezas y el segundo sale pronto.
\end{ejemplo}

\begin{ejemplo}{Ejemplo resuelto 21 --- Neum\'aticos \origen{Serie 3, problema 9}}
$\fxy=K(x^2+y^2)$, $K=\tfrac{3}{380000}$, en $[20,30]\times[20,30]$, con $f_X(x)=\tfrac{3x^2}{38000}+\tfrac1{20}$
(y lo mismo para $Y$ por simetr\'ia).
{\small\begin{align*}
\E(X)&=\int_{20}^{30}x\Big(\tfrac{3x^2}{38000}+\tfrac1{20}\Big)dx=\tfrac{1925}{76}=25.3289=\E(Y)\\
\E(X^2)&=\int_{20}^{30}x^2\Big(\tfrac{3x^2}{38000}+\tfrac1{20}\Big)dx=\tfrac{37040}{57}=649.8246\ \Rightarrow\ \Var(X)=\Var(Y)=8.2690\\
\E(XY)&=K\int_{20}^{30}\!\!\int_{20}^{30}xy(x^2+y^2)\,dx\,dy=\tfrac{24375}{38}=641.4474
\end{align*}}
\textbf{a)} $\Cov(X,Y)=641.4474-(25.3289)^2=\resp{-0.1082}$ (la respuesta $-0.108$ de la serie).

\textbf{b)} $\rho=\dfrac{-0.1082}{8.2690}=\resp{-0.0131}$. Relaci\'on lineal pr\'acticamente nula.

\textbf{Cuidado:} aqu\'i se restan dos n\'umeros casi iguales. Trabaja con fracciones o con al menos 4 decimales; si redondeas
$\E(X)$ a 25.33 la covariancia sale mal.
\end{ejemplo}

\begin{ejemplo}{Ejemplo resuelto 22 \origen{Apuntes, Ej.~3.12 (en blanco), incisos c y d}}
$\fxy=4xye^{-(x^2+y^2)}$. Ya vimos que son independientes con $f_X=2xe^{-x^2}$.

\textbf{c)} Por independencia las esperanzas condicionales son constantes:
$\E[X\mid y]=\E(X)=\int_0^\infty2x^2e^{-x^2}dx=\resp{\tfrac{\sqrt\pi}{2}}$, y $\E[Y\mid x]=\resp{\tfrac{\sqrt\pi}{2}}$.

\textbf{d)} Independientes $\Rightarrow\Cov(X,Y)=0\Rightarrow\resp{\rho=0}$. No hace falta integrar.
\end{ejemplo}

\begin{mecaniza}
\textbf{1.} (Autoexamen 7) $\fxy=4xy$ en el cuadrado unitario. $\Cov(X,Y)$.
\quad\textbf{R:} se factoriza como $(2x)(2y)$ en un rect\'angulo: independientes, $\Cov=0$ (inciso C).

\textbf{2.} (Autoexamen 10) $f(0,1)=0.1$, $f(1,1)=0.3$, $f(0,2)=0.4$, $f(1,2)=0.2$. Halla $\rho$.
\quad\textbf{R:} $\E(X)=0.5$, $\E(Y)=1.6$, $\E(XY)=0.3+0.4=0.7$, $\Cov=-0.1$, $\Var(X)=0.25$, $\Var(Y)=0.24$,
$\rho=-0.1/\sqrt{0.06}=-0.408$ (inciso A).

\textbf{3.} Con $\fxy=24xy$ del tri\'angulo $x+y\le1$: $\Cov$ y $\rho$.
\quad\textbf{R:} $\E(X)=\E(Y)=\tfrac25$, $\E(X^2)=\tfrac15$, $\Var=\tfrac1{25}$, $\E(XY)=\tfrac2{15}$,
$\Cov=\tfrac2{15}-\tfrac4{25}=-\tfrac2{75}$, $\rho=-\tfrac23$.

\textbf{4.} Con $\fxy=x+y$ del Ej.~3.4: $\Cov$ y $\rho$.
\quad\textbf{R:} $\E(X)=\E(Y)=\tfrac7{12}$, $\E(XY)=\tfrac13$, $\Cov=-\tfrac1{144}$, $\Var=\tfrac{11}{144}$, $\rho=-\tfrac1{11}$.
\end{mecaniza}

\section{Combinaciones lineales}

\begin{teoria}{Las dos f\'ormulas}
Para $W=aX+bY+c$:
\[
\E(W)=a\,\E(X)+b\,\E(Y)+c
\qquad
\Var(W)=a^2\Var(X)+b^2\Var(Y)+2ab\,\Cov(X,Y)
\]
\begin{itemize}
\item La constante $c$ \textbf{no cambia} la variancia.
\item Los coeficientes entran \textbf{al cuadrado} en la variancia: un signo menos desaparece ah\'i, pero \textbf{no} en el t\'ermino $2ab\Cov$.
\item Si son independientes, $\Cov=0$ y queda $\Var(W)=a^2\Var(X)+b^2\Var(Y)$.
\item Si te dan $\rho$ en vez de la covariancia: $\Cov=\rho\,\sigma_X\sigma_Y$.
\end{itemize}
Con $n$ variables: $\Var\big(\sum a_iX_i\big)=\sum a_i^2\Var(X_i)+2\sum_{i<j}a_ia_j\Cov(X_i,X_j)$, o en forma matricial
$\mathbf a^{T}\Sigma\,\mathbf a$ con $\Sigma$ la matriz de variancia-covariancia.
\end{teoria}

\begin{ejemplo}{Ejemplo resuelto 23 \origen{Serie 3, problema 11}}
$X$, $Y$ independientes, $\E(X)=3$, $\E(X^2)=25$, $\E(Y)=10$, $\E(Y^2)=164$.

\textbf{a)} $\E(3X+Y-8)=9+10-8=\resp{11}$\qquad
\textbf{b)} $\E(2X-3Y+7)=6-30+7=\resp{-17}$

\textbf{c)} $\Var(X)=25-9=\resp{16}$\quad \textbf{d)} $\sigma_X=\resp{4}$\quad
\textbf{e)} $\Var(Y)=164-100=\resp{64}$\quad \textbf{f)} $\sigma_Y=\resp{8}$

\textbf{g)} $\Var(3X+Y-8)=9(16)+64=\resp{208}$\qquad
\textbf{h)} $\Var(2X-3Y+7)=4(16)+9(64)=\resp{640}$

\textbf{i)} $\E\big[\tfrac{X-3}{4}\big]=\tfrac{3-3}{4}=\resp{0}$,\quad $\Var\big[\tfrac{X-3}{4}\big]=\tfrac{1}{16}(16)=\resp{1}$

\textbf{j)} $\E\big[\tfrac{Y-10}{8}\big]=\resp{0}$,\quad $\Var\big[\tfrac{Y-10}{8}\big]=\tfrac1{64}(64)=\resp{1}$

\textbf{k)} S\'i: para cualquier variable, $Z=\dfrac{X-\mu}{\sigma}$ tiene
$\E(Z)=\dfrac{\mu-\mu}{\sigma}=0$ y $\Var(Z)=\dfrac{\sigma^2}{\sigma^2}=1$ (estandarizaci\'on).
\end{ejemplo}

\begin{ejemplo}{Ejemplo resuelto 24 --- Tres variables \origen{Apuntes, Ej.~3.13 (en blanco)}}
$\E(Y_1)=2$, $\E(Y_2)=-1$, $\E(Y_3)=4$; $\Var=4,\,6,\,8$; $\Cov(Y_1,Y_2)=1$, $\Cov(Y_1,Y_3)=-1$, $\Cov(Y_2,Y_3)=0$.
Para $W=3Y_1+4Y_2-6Y_3$:
\begin{align*}
\E(W)&=3(2)+4(-1)-6(4)=\resp{-22}\\
\Var(W)&=\underbrace{9(4)+16(6)+36(8)}_{420}+2\big[\underbrace{(3)(4)(1)}_{12}+\underbrace{(3)(-6)(-1)}_{18}+\underbrace{(4)(-6)(0)}_{0}\big]=420+60=\resp{480}
\end{align*}
Forma matricial (da lo mismo):
$\begin{bmatrix}3&4&-6\end{bmatrix}
\begin{bmatrix}4&1&-1\\1&6&0\\-1&0&8\end{bmatrix}
\begin{bmatrix}3\\4\\-6\end{bmatrix}=480$.
\end{ejemplo}

\begin{mecaniza}
\textbf{1.} (Autoexamen 8) $\sigma_X^2=3$, $\sigma_Y^2=4$, $\Cov=2$. $\Var(2X-Y)$.
\quad\textbf{R:} $4(3)+1(4)+2(2)(-1)(2)=12+4-8=8$ (inciso E).

\textbf{2.} (Autoexamen 9) $\mu_X=10$, $\sigma_X=15$, $\mu_Y=15$, $\sigma_Y=20$, $\rho=-0.1$; $R=0.4X+0.6Y$.
\quad\textbf{R:} $\E(R)=4+9=13$. $\Cov=-0.1(15)(20)=-30$.
$\Var(R)=0.16(225)+0.36(400)+2(0.4)(0.6)(-30)=36+144-14.4=165.6$, $\sigma_R=12.87$ (inciso D).
\end{mecaniza}

\begin{trampa}
\begin{itemize}
\item Te dan $\sigma$ (desviaci\'on) y la usas como variancia. En la f\'ormula va $\sigma^2$.
\item $\Var(2X-3Y)$ \textbf{no} es $4\Var(X)-9\Var(Y)$: es $4\Var(X)+9\Var(Y)-12\Cov(X,Y)$.
\item Sumar la constante a la variancia. $\Var(X+7)=\Var(X)$.
\item Usar $\E(XY)=\E(X)\E(Y)$ sin que el problema diga (o hayas probado) que son independientes.
\end{itemize}
\end{trampa}

\section{Los errores que m\'as puntos quitan}

\begin{enumerate}
\item \textbf{No dibujar la regi\'on.} Con el dibujo los l\'imites salen solos; sin \'el se pierden el inciso y los que dependen de \'el.
\item \textbf{L\'imites de afuera con variables.} Los de la integral exterior siempre son n\'umeros.
\item \textbf{Olvidar el intervalo} al escribir una marginal o condicional (\guillemotleft$0\le x\le1$; 0 en otro caso\guillemotright).
\item \textbf{Confundir} $P(X<a\mid Y=y_0)$ (se usa la densidad condicional) con $P(X<a\mid Y<b)$ (se usa $P(A\cap B)/P(B)$).
\item \textbf{No sustituir $k$} en los incisos siguientes.
\item \textbf{Afirmar independencia sin mostrarla.} Escribe $f_X$, $f_Y$, su producto y la comparaci\'on.
\item \textbf{Redondear a media operaci\'on} en covariancias. Usa fracciones.
\item \textbf{No resaltar el resultado.} Las instrucciones del examen lo piden: enci\'erralo.
\end{enumerate}

\textbf{Comprobaciones de 10 segundos:} toda tabla, marginal y condicional suma o integra 1; toda probabilidad est\'a
entre 0 y 1; toda variancia es positiva; $-1\le\rho\le1$; $\E(X)$ cae dentro del intervalo de $X$.

\section{Simulacro tipo examen (90 minutos)}

Mismo formato que tu segundo parcial: 5 problemas, 20 puntos cada uno. Resu\'elvelo completo antes de ver las soluciones.

\begin{enumerate}[label=\textbf{\arabic*.-},leftmargin=*]
\item El n\'umero de reclamaciones $X$ y de devoluciones $Y$ que recibe una tienda en un d\'ia tiene la funci\'on de probabilidad conjunta
\begin{center}
\begin{tabular}{c|ccc}
$x\backslash y$ & 0 & 1 & 2\\ \hline
1 & 0.10 & 0.05 & 0.05\\
2 & 0.15 & 0.20 & 0.05\\
3 & 0.05 & 0.15 & 0.20
\end{tabular}
\end{center}
a) Obtener las marginales. b) Calcular $P(X+Y\le3)$. c) Obtener la funci\'on condicional de $Y$ dado $X=2$ y $P(Y\ge1\mid X=2)$.
d) ?`Son independientes? e) Calcular el coeficiente de correlaci\'on.

\item Sea $\fxy(x,y)=k(x+2y)$ para $0\le x\le1$, $0\le y\le2$, y cero en otro caso.
a) Determinar $k$. b) Calcular $P(X+Y\le1)$. c) Obtener las marginales. d) Obtener $f_{Y\mid x}$ y calcular $P(Y<1\mid X=\tfrac12)$.
e) Determinar la curva de regresi\'on de $Y$ dado $x$.

\item Sea $\fxy(x,y)=kx$ para $0\le y\le x\le1$, y cero en otro caso.
a) Determinar $k$. b) Calcular $P(X+Y\le1)$. c) ?`Son independientes? d) Calcular la covariancia y el coeficiente de correlaci\'on.

\item Las variables $X$ y $Y$ tienen $\E(X)=4$, $\Var(X)=9$, $\E(Y)=-2$, $\Var(Y)=4$ y $\rho=-0.5$.
Para $W=2X-3Y+5$ calcular $\E(W)$ y $\Var(W)$. ?`Cu\'anto valdr\'ia $\Var(W)$ si fueran independientes?

\item Un dep\'osito tiene al inicio de la semana una cantidad $Y$ (en miles de litros) de la cual se vende $X$, con
$\fxy(x,y)=6x$ para $0\le x\le y\le1$. Obtener la cantidad promedio que queda en el dep\'osito al final de la semana.
\end{enumerate}

\newpage
\subsection*{Soluciones del simulacro}

\textbf{1.} a) $f_X=(0.20,\,0.40,\,0.40)$ en $x=1,2,3$; $f_Y=(0.30,\,0.40,\,0.30)$ en $y=0,1,2$.

b) Celdas $(1,0),(1,1),(1,2),(2,0),(2,1),(3,0)$: $0.10+0.05+0.05+0.15+0.20+0.05=\resp{0.60}$.

c) Rengl\'on $x=2$ entre $0.40$: $f_{Y\mid X=2}=(0.375,\,0.5,\,0.125)$; $P(Y\ge1\mid X=2)=\resp{0.625}$.

d) $\fxy(1,0)=0.10\ne f_X(1)f_Y(0)=0.06$. \resp{\text{No}}.

e) $\E(X)=2.2$, $\E(Y)=1$, $\E(X^2)=5.4$, $\E(Y^2)=1.6$, $\Var(X)=0.56$, $\Var(Y)=0.6$.
$\E(XY)=0.05+0.10+0.40+0.20+0.45+1.20=2.4$. $\Cov=2.4-2.2=0.2$.
$\rho=0.2/\sqrt{0.56\cdot0.6}=\resp{0.345}$.

\textbf{2.} a) $\int_0^2\!\int_0^1(x+2y)\,dx\,dy=\int_0^2(\tfrac12+2y)\,dy=5\Rightarrow\resp{k=\tfrac15}$.

b) $\tfrac15\int_0^1\!\int_0^{1-x}(x+2y)\,dy\,dx=\tfrac15\int_0^1\big[x(1-x)+(1-x)^2\big]dx=\tfrac15\int_0^1(1-x)\,dx=\resp{\tfrac1{10}}$.

c) $f_X(x)=\tfrac15\int_0^2(x+2y)\,dy=\resp{\tfrac{2x+4}{5}}$, $0\le x\le1$;\quad
$f_Y(y)=\tfrac15\int_0^1(x+2y)\,dx=\resp{\tfrac{1+4y}{10}}$, $0\le y\le2$.

d) $f_{Y\mid x}=\dfrac{(x+2y)/5}{(2x+4)/5}=\resp{\dfrac{x+2y}{2x+4}}$, $0\le y\le2$.
Con $x=\tfrac12$: $\int_0^1\dfrac{0.5+2y}{5}\,dy=\dfrac{0.5+1}{5}=\resp{0.3}$.

e) $\E[Y\mid x]=\int_0^2y\,\dfrac{x+2y}{2x+4}\,dy=\dfrac{2x+\frac{16}{3}}{2x+4}=\resp{\dfrac{3x+8}{3(x+2)}}$.

\textbf{3.} a) $\int_0^1\!\int_0^xkx\,dy\,dx=k\int_0^1x^2dx=\tfrac k3\Rightarrow\resp{k=3}$.

b) La regi\'on $x+y\le1$ dentro de $y\le x$ es el tri\'angulo con v\'ertices $(0,0)$, $(1,0)$, $(\tfrac12,\tfrac12)$. Conviene flecha horizontal:
para $0\le y\le\tfrac12$, entra en $x=y$ y sale en $x=1-y$:
\[\int_0^{1/2}\!\!\int_y^{1-y}3x\,dx\,dy=\tfrac32\int_0^{1/2}\big[(1-y)^2-y^2\big]dy=\tfrac32\int_0^{1/2}(1-2y)\,dy=\tfrac32\cdot\tfrac14=\resp{\tfrac38}\]

c) Soporte triangular: \resp{\text{no}}. Con marginales: $f_X=3x^2$, $f_Y=\tfrac32(1-y^2)$, y $f_Xf_Y\ne3x$.

d) $\E(X)=\int_0^13x^3dx=\tfrac34$, $\E(X^2)=\tfrac35$, $\Var(X)=\tfrac35-\tfrac9{16}=\tfrac3{80}$.
$\E(Y)=\int_0^1\tfrac32y(1-y^2)\,dy=\tfrac38$, $\E(Y^2)=\tfrac15$, $\Var(Y)=\tfrac15-\tfrac9{64}=\tfrac{19}{320}$.
$\E(XY)=\int_0^1\!\int_0^x3x^2y\,dy\,dx=\int_0^1\tfrac32x^4dx=\tfrac3{10}$.
$\Cov=\tfrac3{10}-\tfrac9{32}=\resp{\tfrac3{160}=0.01875}$;\quad
$\rho=\dfrac{3/160}{\sqrt{(3/80)(19/320)}}=\resp{0.397}$.

\textbf{4.} $\E(W)=2(4)-3(-2)+5=\resp{19}$. $\Cov=\rho\,\sigma_X\sigma_Y=-0.5(3)(2)=-3$.
$\Var(W)=4(9)+9(4)+2(2)(-3)(-3)=36+36+36=\resp{108}$. Independientes: $36+36=\resp{72}$.

\textbf{5.} $\E(Y-X)=\int_0^1\!\int_0^y6x(y-x)\,dx\,dy=\int_0^1\big[3x^2y-2x^3\big]_0^ydy=\int_0^1y^3dy=\resp{\tfrac14}$ (250 litros).

\section{Respuestas del autoexamen de tus apuntes}

\begin{center}
\begin{tabular}{c|cccccccccc}
Pregunta & 1 & 2 & 3 & 4 & 5 & 6 & 7 & 8 & 9 & 10\\ \hline
Inciso & B & A & B & B & B & D & C & E & D & A\\
Valor & $\tfrac1{500}$ & 0.03125 & $\tfrac x6-\tfrac1{12}$ & $\tfrac{2x}{y^2}$ & 0.8 & 1.88 & 0 & 8 & 13, 12.87 & $-0.408$
\end{tabular}
\end{center}
El procedimiento de cada una est\'a en los cuadros \guillemotleft Mecaniza\guillemotright{} de las secciones anteriores.

\section{Plan de repaso y hoja de bolsillo}

\textbf{Orden sugerido} (unas 5 horas en total): secciones 2 y 3 (1 h), secci\'on 4 con l\'apiz (1.5 h),
secciones 5 a 7 (1.5 h), simulacro con reloj (1.5 h incluyendo revisi\'on). Lo que falles en el simulacro, rep\'itelo al d\'ia siguiente.

\begin{teoria}{Todo el tema en una caja}
\begin{center}
\begin{tabular}{ll}
Hallar $k$ & $\iint\fxy=1$\\
Probabilidad & sumar celdas / integrar sobre la regi\'on dibujada\\
Marginal & sumar o integrar la \emph{otra} variable\\
Condicional & conjunta $\div$ marginal de lo conocido\\
Independencia & ?`$\fxy=f_Xf_Y$ en todo punto? (soporte no rectangular: no)\\
$\E[g(X,Y)]$ & $\sum\sum g\,\fxy$ \ o \ $\iint g\,\fxy$\\
Regresi\'on & $\E[Y\mid x]=\int y\,f_{Y\mid x}\,dy$\\
Covariancia & $\E(XY)-\E(X)\E(Y)$\\
Correlaci\'on & $\rho=\Cov/(\sigma_X\sigma_Y)$, \ $\rho^2$ = determinaci\'on\\
Combinaci\'on lineal & $\E=a\mu_X+b\mu_Y+c$; \ $\Var=a^2\sigma_X^2+b^2\sigma_Y^2+2ab\Cov$
\end{tabular}
\end{center}
\end{teoria}

\textbf{Calculadora TI-Nspire CX II.} Esta versi\'on no incluye todav\'ia las indicaciones de qu\'e pesta\~na del programa usar en cada
tipo de ejercicio. Se agregar\'an cuando el programa est\'e hecho, como acordamos.


\end{document}
